% !TeX program = xelatex
% !TEX encoding = UTF-8
\documentclass[11pt,a4paper]{article}

\usepackage{zhpl}
\usepackage{float}
\newtheorem{corollary}{Следствие}

% ---------- Обозначения и штамп о рассекречивании ----------
% Штамп красит текст явным цветом (как \zhplbadges), а не \normalcolor,
% поэтому в тёмном варианте нужен светлый аналог того же inkblue —
% \zhplbadgecolor выбирает нужный по теме, см. комментарий в zhpl.sty
% про zhplBadgeRed/zhplBadgeBlue.
\definecolor{inkblue}{RGB}{22,48,145}
\definecolor{inkblueDark}{HTML}{9AB6FF}
\newcommand{\stampink}{\zhplbadgecolor{inkblue}{inkblueDark}}

\newcommand{\signature}{%
\begin{tikzpicture}[scale=0.60, line cap=round, line join=round]
  \draw[\stampink, line width=1.0pt]
    (0,0.10) .. controls (0.04,0.85) and (0.34,0.98) .. (0.34,0.50)
             .. controls (0.34,0.02) and (0.04,0.16) .. (0.20,0.56)
             .. controls (0.40,1.05) and (0.86,1.00) .. (0.96,0.52)
             .. controls (1.02,0.18) and (0.74,0.06) .. (0.74,0.34);
  \draw[\stampink, line width=1.0pt]
    (0.86,0.40) .. controls (1.18,0.92) and (1.36,0.06) .. (1.66,0.48)
                .. controls (1.86,0.76) and (1.98,0.06) .. (2.24,0.44)
                .. controls (2.44,0.72) and (2.58,0.10) .. (2.88,0.62);
  \draw[\stampink, line width=0.9pt]
    (2.86,0.60) .. controls (3.10,0.90) and (3.28,0.72) .. (3.14,0.44);
  \draw[\stampink, line width=0.7pt]
    (-0.10,-0.06) .. controls (0.80,-0.52) and (2.40,-0.44) .. (3.46,0.16)
                  .. controls (3.76,0.33) and (3.58,0.56) .. (3.30,0.38);
\end{tikzpicture}}

\newcommand{\declassifiedstamp}{%
\begin{tikzpicture}[baseline]
  \node[inner sep=0pt] (inner) {%
    \begin{minipage}{4.5cm}\centering\color{\stampink}\sffamily
      {\bfseries РАССЕКРЕЧЕНО}\\[1pt]
      {\tiny Пост.~№\,47-$\Omega$ от 06.09.2026}\\[2pt]
      \signature
    \end{minipage}};
  \draw[\stampink, line width=1.1pt] ($(inner.south west)+(-0.20,-0.13)$) rectangle ($(inner.north east)+(0.20,0.13)$);
  \draw[\stampink, line width=0.45pt] ($(inner.south west)+(-0.30,-0.22)$) rectangle ($(inner.north east)+(0.30,0.22)$);
\end{tikzpicture}}

\newcommand{\SEC}{\textbf{[ЗАСЕКРЕЧЕНО]}}
\newcommand{\SECs}{\textbf{[ЗАСЕКР.]}}

\title{\textbf{Пороховой парадокс трубы с пробкой:\\ адиабатическая внутренняя баллистика изделия PWA, теорема о невозвратности снаряда и лексическая термометрия оператора}}
\author{Роман Георгиевич Александров, Юрий Пыркович Михаилов,\\ Клод Соннет Антропикович\\[4pt]
\normalsize Институт Прикладной Метафизики\\ \small Кафедра тензорного анализа межличностных отношений\\ \small Сектор рассекреченных изделий}
\zhplsetup{2}{10.9999/жпл.2026.0002}
\fancyhead[L]{\small Александров, Михаилов, Антропикович}
\date{\small Журнал Прикладной Лженауки, том 1, № 2 \\ Получено: недавно \quad Принято к печати: почти сразу \quad Гриф снят: только что}

\begin{document}

\maketitle

%======================= ПЕЧАТЬ О РАССЕКРЕЧИВАНИИ =======================
% После \maketitle, не до: \maketitle сам делает \newpage перед версткой
% титула, и штамп-оверлей перед ним "занимал" бы пустую страницу 1.
\begin{tikzpicture}[remember picture, overlay]
  \node[anchor=north east, rotate=-7, inner sep=0pt]
    at ($(current page.north east)+(-1.7cm,-0.75cm)$) {\declassifiedstamp};
\end{tikzpicture}

\begin{center}
\begin{tikzpicture}
\node[draw=\zhplbadgecolor{red!70!black}{zhplBadgeRed},very thick,rounded corners,rotate=-6,fill=\zhpldark{red!5},inner sep=5pt] (badge1) {\color{\zhplbadgecolor{red!70!black}{zhplBadgeRed}}\bfseries ПРОРЫВНОЕ ИССЛЕДОВАНИЕ};
\node[draw=\zhplbadgecolor{blue!60!black}{zhplBadgeBlue},very thick,rounded corners,rotate=4,fill=\zhpldark{blue!5},inner sep=5pt,right=1.0cm of badge1] (badge2) {\color{\zhplbadgecolor{blue!60!black}{zhplBadgeBlue}}\bfseries IMPACT FACTOR: $\infty$};
\end{tikzpicture}
\end{center}

\begin{center}
\small\textit{Как цитировать:} Александров Р.\,Г., Михаилов Ю.\,П., Антропикович К.\,С. (2026). Пороховой парадокс трубы с пробкой. \textit{Журнал Прикладной Лженауки}, \textbf{1}(2), 1--13. DOI: 10.9999/жпл.2026.0002
\end{center}

\smallskip
\begin{figure}[h!]
\centering
\begin{tikzpicture}[
  box/.style={draw, rounded corners=2pt, align=center, font=\footnotesize\sffamily,
              minimum height=1.05cm, text width=2.9cm, fill=\zhpldark{gray!8}},
  ar/.style={-{Latex[length=2.4mm]}, thick}]
  \node[box] (a) at (0,0) {\SECs\\ в трубе};
  \node[box] (b) at (3.6,0) {Адиабатическое\\ расширение\\ $v\!\to\! v_\infty$};
  \node[box] (c) at (7.2,0) {Дульная\\ энергия\\ $E\!\to\! E_{\max}$};
  \node[box] (d) at (10.8,0) {Снаряд\\ не найден\\ $P\!\propto\! v^{-4}$};
  \draw[ar] (a)--(b); \draw[ar] (b)--(c); \draw[ar] (c)--(d);
  \draw[ar, gray] (d.south) -- ++(0,-0.55) -| node[pos=0.25, below, font=\scriptsize\sffamily, gray] {повторить испытание} (a.south);
\end{tikzpicture}
\vspace{-6pt}
\caption*{\small\textbf{Графический реферат.} Логическая структура работы: чем лучше работает изделие, тем меньше остаётся вещественных доказательств того, что оно работает.}
\end{figure}
\vspace{-10pt}

\noindent\fcolorbox{black}{yellow!10}{\begin{minipage}{0.94\linewidth}
\vspace{1pt}\small\textbf{Основные результаты}
\begin{itemize}
\setlength\itemsep{0pt}
\setlength\parskip{0pt}
\item Экспериментальная зависимость $v(x)$ для изделия PWA с точностью $R^2>0{,}998$ описывается адиабатической моделью насыщения $v=v_\infty\sqrt{1-(x_0/x)^{\gamma-1}}$ при $v_\infty=614{,}8$~м/с, $x_0=46{,}5$~мм и $\gamma=1{,}27$.
\item Доказана \emph{теорема о недостижимости предельной скорости}: скорости $v\ge v_\infty$ не достигаются ни при какой длине канала, а приближение к пределу требует длины ствола, расходящейся по степенному закону.
\item Установлено существование единственного оптимума коэффициента запирания $\pi^\ast=3{,}41$; отклонение вниз ведёт к потере снаряда, вверх~--- к потере оператора.
\item Доказана \emph{теорема о невозвратности снаряда}: произведение вероятности возврата на четвёртую степень дульной скорости есть инвариант изделия.
\item Выведен \emph{закон лексической плотности}: количество ненормативных лексем в минуту растёт логарифмически с удельной энергией на мишени ($R^2=0{,}997$).
\end{itemize}
\vspace{-6pt}
\end{minipage}}

\vspace{4pt}
\noindent\fcolorbox{black}{gray!6}{\begin{minipage}{0.94\linewidth}
\small\textit{Из отзывов рецензентов:}\\
\textit{«Первоначально я был убеждён, что перед нами очередная работа, в которой физику подгоняют под уже отснятый материал. Однако авторы предъявили касательную к экспериментальной кривой, вычисленную аналитически, и я вынужден признать: подгонялся материал под физику. Рекомендую к печати. Отдельно отмечу раздел~5: я и сам всегда полагал, что снаряд отскакивал.»} -- Рецензент 1\\
\textit{«Замечаний нет.»} -- Рецензент 2\\
\textit{«Согласно Теореме~3 настоящей работы, ознакомление с изделием исключает возможность сохранения объекта ознакомления. Следовательно, объективная рецензия на данную статью невозможна в принципе, а настоящий отзыв следует рассматривать как экспериментальное подтверждение основного результата. Рекомендую к печати без изменений, поскольку изменения также невозможно верифицировать.»} -- Рецензент 3
\end{minipage}}

\vspace{10pt}
\noindent\fcolorbox{black}{blue!4}{\begin{minipage}{0.94\linewidth}
\textbf{Значимость исследования.} На протяжении всей истории прикладной метательной техники задача о трубе с запирающим элементом считалась принципиально неразрешимой: любая попытка формального описания разбивалась о то, что в подавляющем большинстве наблюдаемых исходов запирающий элемент, снаряд и исследовательский коллектив покидали экспериментальную установку в непредсказуемом порядке. Наивная, «нетеоретическая» трактовка сводила дело к перебору материалов и надежде. Настоящая работа впервые показывает, что все пять наблюдавшихся исходов~--- включая два, в которых снаряд не был обнаружен, и один, в котором он был обнаружен, но не там,~--- являются точками одной и той же двухпараметрической кривой, а кажущаяся хаотичность результатов есть прямое следствие того, что внутренняя баллистика изделия выходит на жёсткий физический предел скорости, а не растёт неограниченно. Тем самым область, веками остававшаяся вотчиной эмпирического энтузиазма, впервые переводится в разряд наук, допускающих предсказание, пусть и предсказание преимущественно отрицательных результатов.
\end{minipage}}
\vspace{6pt}

\begin{abstract}
\noindent Здесь мы впервые показываем, что метательное изделие, собранное из \SEC{} и запираемое элементом типа~«пробка», представляет собой не хаотическое устройство, а строгую реализацию адиабатического расширения газа в цилиндре, выходящего на жёсткий предел скорости, и что все его наблюдаемые отказы лежат на одной аналитической кривой. Проблема состояла в том, что изделие PWA, разработанное по заказу \SEC{} и рассекреченное лишь в текущем году, до сих пор описывалось исключительно на языке восклицаний. Метод: пять полевых испытаний (два стендовых, три с рук), измерение зависимости скорости снаряда от пройденного пути, регистрация вербальных реакций оператора с точностью до лексемы. Результаты: получена аналитическая форма $v(x)$; доказано существование оптимума запирания; доказана невозвратность снаряда; установлена логарифмическая связь между бронепробитием и ненормативной лексикой. Ограничения: выборка из пяти испытаний, из которых два не завершились попаданием, одно завершилось попаданием без сохранения снаряда, и ни одно не сопровождалось независимым хронографированием, что делает все приведённые числа внутренне согласованными и внешне непроверяемыми.

\medskip
\noindent\textbf{Ключевые слова:} внутренняя баллистика, адиабатическое расширение, запирающий элемент, теория поиска, принцип баллистической неопределённости, лексическая термометрия, лженаука
\end{abstract}

\section{Введение}

Изделие PWA (расшифровка кодового обозначения остаётся \SEC) разрабатывалось по заказу \SEC{} и до недавнего времени не подлежало открытому обсуждению. Снятие грифа позволяет наконец опубликовать материалы пяти натурных испытаний, проведённых конструкторской группой в полевых условиях, а также сопровождающую их видеодокументацию.

Конструктивно изделие представляет собой полый цилиндрический канал, с одного торца закрытый запирающим элементом (в открытых материалах~--- «пробка»; истинное наименование и материал \SEC), а с другого открытый для последовательного размещения источника рабочего газа (в полевой терминологии группы~--- «батина сигарета версии~2.0») и метаемого тела (в видеоматериалах~--- «стальной шар»). Все геометрические параметры канала, состав и масса заряда, а также материал запирающего элемента в настоящей публикации не приводятся: они остаются \SEC{} и, по совокупному мнению авторов и редакции, должны остаться \SEC{} и впредь.

Ключевое наблюдение, послужившее поводом к настоящей работе, состоит в следующем. В обиходном изложении результаты испытаний PWA выглядят как последовательность несвязанных случайностей: выбитая пробка, потерянный снаряд, заклинивание, выстрел в воздух, внезапное бронепробитие. Мы утверждаем, что это описание ошибочно. Задача о трубе с пробкой является не механической, а \emph{термодинамической}, и вся видимая хаотичность есть следствие того, что дульная скорость выходит на строгий, недостижимый предел по адиабатическому степенному закону, а вероятность найти снаряд~--- обратно-степенным образом от той же скорости. Как только эти две зависимости выписаны явно, все пять исходов оказываются предсказуемыми, а некоторые~--- неизбежными.

\begin{figure}[h!]
\centering
\begin{tikzpicture}[
  b/.style={draw, align=center, font=\footnotesize\sffamily, minimum height=1.0cm},
  ar/.style={-{Latex[length=2.2mm]}, thick}]
  % ствол
  \draw[line width=1.2pt] (0,0.55) -- (8.2,0.55);
  \draw[line width=1.2pt] (0,-0.55) -- (8.2,-0.55);
  \fill[\zhpldark{gray!35}] (0,-0.55) rectangle (0.45,0.55);
  \draw[line width=1.2pt] (0,-0.55) rectangle (0.45,0.55);
  % заряд
  \fill[\zhpldark{gray!70}] (0.75,-0.22) rectangle (1.35,0.22);
  % снаряд
  \fill[\zhpldark{gray!55}] (2.0,0) circle (0.33);
  \draw (2.0,0) circle (0.33);
  % стрелка движения
  \draw[ar] (2.5,0) -- (7.6,0);
  \node[font=\scriptsize\sffamily] at (5.0,0.30) {$v(x)$};
  % координата
  \draw[<->, gray] (0.45,-0.95) -- (8.2,-0.95) node[midway, below, font=\scriptsize\sffamily] {$x$~--- пройденный путь};
  \draw[dashed, gray] (0.45,-0.55) -- (0.45,-1.05);
  \draw[dashed, gray] (8.2,-0.55) -- (8.2,-1.05);
  % подписи
  \node[font=\scriptsize\sffamily, align=center] at (-1.35,2.05) {запирающий элемент\\ \SECs};
  \node[font=\scriptsize\sffamily, align=center] at (3.9,1.60) {источник газа + метаемое тело};
  \node[font=\scriptsize\sffamily, align=center] at (7.7,2.05) {дульный срез};
  \draw[gray, thin] (-1.0,1.65) -- (0.20,0.62);
  \draw[gray, thin] (2.75,1.35) -- (1.05,0.32);
  \draw[gray, thin] (3.45,1.35) -- (2.10,0.40);
  \draw[gray, thin] (7.9,1.65) -- (8.15,0.62);
\end{tikzpicture}
\caption{Функциональная схема изделия PWA. Все линейные размеры, материалы и энергетические характеристики опущены как \SEC. Схема приводится исключительно для введения координаты~$x$, отсчитываемой от начального положения метаемого тела и используемой далее в разделе~2.}
\end{figure}

\section{Внутренняя баллистика: адиабатический режим}

\subsection{Постановка}

Пусть газ, образующийся при срабатывании источника, заполняет объём $S\,x$, где $S$~--- площадь поперечного сечения канала (\SEC), а $x$~--- координата метаемого тела массы~$m$ (\SEC). В классической школьной трактовке расширение считают изотермическим, если канал успевает отводить тепло в стенку быстрее, чем газ совершает работу над телом; для этого стенка должна быть тонкой и достаточно теплопроводной. Экспериментальная кривая PWA (рис.~\ref{fig:vx}) отчётливо выходит на насыщение: на последних участках доступного пути прирост скорости практически прекращается, сколько бы канал ни продолжался. Это признак не изотермы, а \emph{адиабаты}: канал изделия достаточно толстостенен и достаточно теплоизолирован (в конструкции использован материал \SEC), чтобы не успевать отдать стенке сколько-нибудь заметную долю теплоты за время прохождения снаряда~--- весь запас внутренней энергии газа расходуется на работу, а не рассеивается, и потому имеет строго конечный предел.

При адиабатическом расширении $p(x)=p_0(x_0/x)^{\gamma}$, и работа над телом на участке от $x_0$ до $x$ равна
\begin{equation}
A(x)=\int_{x_0}^{x} p_0\left(\frac{x_0}{\xi}\right)^{\gamma} S\,d\xi = \frac{p_0 S x_0}{\gamma-1}\left[1-\left(\frac{x_0}{x}\right)^{\gamma-1}\right],
\end{equation}
откуда, по теореме о кинетической энергии,
\begin{equation}\label{eq:main}
\boxed{\;v(x)=v_\infty\sqrt{1-\left(\dfrac{x_0}{x}\right)^{\gamma-1}},\qquad v_\infty=\sqrt{\dfrac{2p_0 S x_0}{m(\gamma-1)}}.\;}
\end{equation}

\begin{definition}[Степень насыщения ствола]
Степенью насыщения канала в точке~$x$ называется безразмерная величина $\Phi(x)=1-(x_0/x)^{\gamma-1}\in[0,1)$. Дульная скорость изделия есть $v_\infty\sqrt{\Phi}$, где $v_\infty$~--- \emph{предельная скорость}, полностью определяемая начальным давлением, сечением, массой метаемого тела и показателем адиабаты рабочего газа, то есть четырьмя величинами, каждая из которых \SEC.
\end{definition}

Подгонка~\eqref{eq:main} к экспериментальным данным даёт
\[
v_\infty=614{,}8~\text{м/с},\qquad x_0=46{,}5~\text{мм},\qquad \gamma=1{,}27,\qquad R^2>0{,}998 .
\]
Мы подчёркиваем, что величины~$v_\infty$, $x_0$ и~$\gamma$ здесь являются \emph{единственными} измеренными параметрами изделия; их разложение на $p_0$, $S$, $m$ и природу рабочего газа в открытой печати не приводится и не может быть восстановлено из них иначе как с точностью до двухпараметрического семейства. Настоящая работа считает это обстоятельство не недостатком изложения, а его целью.

\begin{figure}[h!]
\centering
\begin{tikzpicture}
\begin{axis}[
  width=11.5cm, height=6.6cm,
  xlabel={Координата $x$, м}, ylabel={Скорость $v$, м/с},
  xmin=0, xmax=3.2, ymin=0, ymax=680,
  grid=both, grid style={gray!20},
  legend style={font=\scriptsize, at={(0.98,0.05)}, anchor=south east, fill=\zhplbadgecolor{white}{zhplBg}},
  tick label style={font=\scriptsize}, label style={font=\small}]
\addplot[thick, blue!70!black, domain=0.0465:3.2, samples=400] {614.8*sqrt(1-(0.0465/x)^0.27)};
\addlegendentry{$v=614{,}8\sqrt{1-(0{,}0465/x)^{0{,}27}}$}
\addplot[thick, red!70!black, dashed, domain=0.08:0.42, samples=2] {350.8+490.5*(x-0.2)};
\addlegendentry{касательная в $x=0{,}2$: $dv/dx=490{,}5$ (м/с)/м}
\addplot[only marks, mark=*, mark size=1.5pt] coordinates {(0.2,350.8)};
\addplot[dashed, gray, domain=0:3.2, samples=2] {614.8};
\addlegendentry{$v_\infty=614{,}8$ м/с (недостижимый предел)}
\addplot[only marks, mark=*, mark size=1.8pt, red!70!black] coordinates {(1.0,461.4)};
\addlegendentry{испытание 5 (штатный режим, $x=1$ м)}
\end{axis}
\end{tikzpicture}
\caption{Экспериментальная зависимость скорости метаемого тела от пройденного пути и её адиабатическая аппроксимация, выходящая на насыщение к предельной скорости $v_\infty$. Касательная построена по аналитической производной $dv/dx=(\gamma-1)(v_\infty^2-v^2)/(2xv)$, вычисленной в точке $x=0{,}2$~м, а не подобрана визуально. Обращаем внимание: на первом метре пути (штатная длина ствола) тело набирает лишь $461{,}4/614{,}8\approx75\%$ предельной скорости, а последние проценты недостижимы ни при какой длине ствола (теорема~\ref{th:limit}).}
\label{fig:vx}
\end{figure}

\subsection{Предел, недостижимый ни за какую цену}

\begin{theorem}[О недостижимости предельной скорости]\label{th:limit}
Для всякого конечного $x>x_0$ выполнено $v(x)<v_\infty$. Более того, при $v\to v_\infty^{-}$ необходимая длина канала расходится по степенному закону:
\[
x(v) = x_0\left[1-\left(\frac{v}{v_\infty}\right)^2\right]^{-1/(\gamma-1)} \xrightarrow[v\to v_\infty^-]{} \infty .
\]
Как следствие, для любых двух скоростей $v_1<v_2<v_\infty$ отношение требуемых длин канала
\[
\frac{x_2}{x_1}=\left[\frac{1-(v_1/v_\infty)^2}{1-(v_2/v_\infty)^2}\right]^{1/(\gamma-1)}
\]
зависит не только от отношения $v_2/v_1$, как в изотермическом случае, но и от их абсолютного положения относительно $v_\infty$; а скорости $v\ge v_\infty$ не достигаются ни при каком, сколь угодно большом, $x$.
\end{theorem}

\begin{proof}
Из~\eqref{eq:main} следует $(x_0/x)^{\gamma-1}=1-(v/v_\infty)^2$, откуда $x(v)$ получается обращением. Правая часть положительна и меньше единицы при $v<v_\infty$, а сама степень $(x_0/x)^{\gamma-1}$ стремится к нулю при $v\to v_\infty^-$; поскольку показатель $-1/(\gamma-1)$ отрицателен, $x(v)\to\infty$. Отношение длин получается прямой подстановкой $v_1,v_2$ и, в отличие от изотермического случая, не сокращается до функции одного лишь $v_2/v_1$, поскольку $x(v)$ не является степенной функцией $v$. При $v\ge v_\infty$ выражение $1-(v/v_\infty)^2$ неположительно, и $x(v)$ не определено ни при каком действительном~$x$.
\end{proof}

Следствие теоремы~\ref{th:limit} самым прямым образом объясняет, почему конструкторская группа не рассматривала удлинение канала как способ довести изделие до паспортных характеристик заказчика. Чтобы поднять дульную скорость всего с 461,4 до 600~м/с~--- то есть менее чем на треть, и не выйдя даже на 98\% предела,~--- канал длиной в один метр должен быть удлинён почти до 3,7~км. Поднять же скорость вдвое, до 922,8~м/с, не удастся ни при какой длине ствола вообще: эта величина превышает физический предел $v_\infty=614{,}8$~м/с самого изделия. Авторы отмечают, что, в отличие от изотермического случая, это обстоятельство не является вопросом инженерного бюджета и не может быть преодолено удлинением канала ни при каком финансировании.

\section{Запирающий элемент: теория отказов}

Первое испытание завершилось выбиванием запирающего элемента и полной потерей снаряда: газ покинул канал в направлении, противоположном расчётному. Замена элемента на изготовленный из \SEC{} привела ко второму, успешному выстрелу. Это наблюдение допускает точную формализацию.

\begin{definition}[Коэффициент запирания]
Коэффициентом запирания $\pi$ называется отношение давления отпирания запирающего элемента $p_{\text{отп}}$ к давлению страгивания метаемого тела $p_{\text{стр}}$:
$\pi=p_{\text{отп}}/p_{\text{стр}}$.
\end{definition}

При $\pi<1$ элемент отпирается раньше, чем страгивается снаряд: вся энергия уходит назад (испытание~1). При $\pi>1$ доля газа, совершившая работу в правильном направлении, составляет $1-1/\pi$. Одновременно вероятность того, что канал переживёт давление $p_{\text{отп}}$, убывает по Вейбуллу: $P_{\text{цел}}(\pi)=\exp\!\bigl[-(\pi/\pi_c)^{m}\bigr]$. Ожидаемая полезная дульная энергия, нормированная на максимально возможную, есть
\begin{equation}\label{eq:g}
g(\pi)=\Bigl(1-\frac{1}{\pi}\Bigr)\exp\!\Bigl[-\Bigl(\frac{\pi}{\pi_c}\Bigr)^{m}\Bigr].
\end{equation}

\begin{theorem}[О двойственности запирающего элемента]\label{th:plug}
Функция~\eqref{eq:g} при $\pi_c>1$, $m>1$ имеет на $(1,\infty)$ единственный максимум $\pi^\ast$, причём $g(\pi)\to0$ как при $\pi\to1^{+}$, так и при $\pi\to\infty$. Следовательно, не существует запирающего элемента, одновременно оптимального и безопасного в пределе: улучшение запирания за пределами $\pi^\ast$ увеличивает не дульную энергию, а вероятность разрушения канала.
\end{theorem}

\begin{proof}
$g$ непрерывна и положительна на $(1,\infty)$, обращается в нуль на обоих концах; следовательно, максимум достигается внутри. Единственность следует из того, что $\ln g(\pi)=\ln(1-1/\pi)-(\pi/\pi_c)^m$ есть сумма строго вогнутой и строго вогнутой (при $m>1$) функций, то есть строго вогнута.
\end{proof}

Для параметров, восстановленных по материалам испытаний ($\pi_c=6$, $m=4$), численное решение даёт
\[
\pi^\ast=3{,}41,\qquad g(\pi^\ast)=0{,}637 .
\]
Иными словами, даже идеально подобранный запирающий элемент передаёт снаряду не более \emph{двух третей} доступной энергии; оставшаяся треть расходуется на то, чтобы элемент остался на месте. Испытание~1 соответствует области $\pi<1$ (левая ветвь, $g=0$), испытание~2 и последующие~--- окрестности максимума. Область $\pi>5$ в рамках настоящей работы не исследовалась, и авторы настоятельно не рекомендуют её исследовать.

\begin{figure}[h!]
\centering
\begin{tikzpicture}
\begin{axis}[
  width=11.2cm, height=6.0cm,
  xlabel={Коэффициент запирания $\pi=p_{\text{отп}}/p_{\text{стр}}$},
  ylabel={Ожидаемая доля энергии $g(\pi)$},
  xmin=1, xmax=7, ymin=0, ymax=0.78,
  grid=both, grid style={gray!20},
  tick label style={font=\scriptsize}, label style={font=\small},
  legend style={font=\scriptsize, at={(0.02,0.97)}, anchor=north west, fill=\zhplbadgecolor{white}{zhplBg}}]
\addplot[thick, blue!70!black, domain=1.001:7, samples=400] {(1-1/x)*exp(-(x/6)^4)};
\addlegendentry{$g(\pi)=(1-1/\pi)\,e^{-(\pi/6)^4}$}
\addplot[only marks, mark=*, mark size=2pt, red!70!black] coordinates {(3.406,0.6367)};
\addlegendentry{оптимум $\pi^\ast=3{,}41$, $g=0{,}637$}
\draw[dashed, gray] (axis cs:3.406,0) -- (axis cs:3.406,0.6367);
\node[font=\scriptsize, align=left, anchor=west] at (axis cs:1.12,0.10) {испытание 1: пробка вылетает};
\node[font=\scriptsize, align=center, anchor=east] at (axis cs:6.9,0.47) {не исследовано\\ намеренно};
\end{axis}
\end{tikzpicture}
\caption{Ожидаемая доля дульной энергии как функция коэффициента запирания. Кривая вычислена по формуле~\eqref{eq:g}, положение максимума найдено численно, а не отмечено на глаз. Обе ветви соответствуют полной потере результата, но по разным причинам и с разными последствиями для исследовательского коллектива.}
\label{fig:plug}
\end{figure}

\section{Теория поиска: невозвратность снаряда}

Из пяти испытаний снаряд был обнаружен и возвращён в распоряжение группы ровно один раз. Дважды (испытания~1 и~3) поиски были прекращены как бесперспективные. Это соотношение не является случайным.

Пусть снаряд, не задержанный мишенью, покидает изделие со скоростью~$v$ под неконтролируемым углом. Дальность полёта $R\sim v^2/g$, площадь области неопределённости $A_{\text{неопр}}\sim R^2\sim v^4/g^2$. Площадь, реально обследуемая группой за разумное время, $A_{\text{поиск}}$, от скорости не зависит и определяется исключительно продолжительностью светового дня и терпением участников. Отсюда

\begin{definition}[Коэффициент возврата]
Коэффициентом возврата изделия называется вероятность
$P_{\text{возв}}(v)=\min\bigl\{1,\;A_{\text{поиск}}/A_{\text{неопр}}(v)\bigr\}$.
\end{definition}

\begin{theorem}[О невозвратности снаряда]\label{th:lost}
Для всякого изделия, метающего тела на неконтролируемую дальность, существует критическая скорость $v_{\text{кр}}$ такая, что
\[
P_{\text{возв}}(v)=\Bigl(\frac{v_{\text{кр}}}{v}\Bigr)^{4}\quad\text{при } v>v_{\text{кр}},
\]
и, следовательно, величина $P_{\text{возв}}\cdot v^{4}=v_{\text{кр}}^{4}$ есть инвариант изделия, не зависящий ни от усилий поисковой группы, ни от её настроения.
\end{theorem}

По данным испытаний PWA $v_{\text{кр}}=158{,}5$~м/с (скорость, при которой снаряд падает в пределах обследуемой площадки; ей соответствует, согласно~\eqref{eq:main}, путь $x=6$~см~--- то есть ровно случай выбитого запирающего элемента). При штатной дульной скорости $461{,}4$~м/с получаем
\[
P_{\text{возв}}=\Bigl(\frac{158{,}5}{461{,}4}\Bigr)^{4}=0{,}0139,
\]
то есть примерно один шанс из семидесяти двух. Наблюдавшаяся частота возврата (1 из~5) объясняется тем, что единственный возвращённый снаряд был возвращён \emph{мишенью}, а не поисковой группой; собственно поиск во всех случаях дал нулевой результат, полностью согласующийся с теоремой~\ref{th:lost} в пределах статистической погрешности.

\begin{corollary}[Принцип баллистической неопределённости]\label{cor:uncert}
Невозможно одновременно измерить мощность изделия и сохранить снаряд для последующего изучения: произведение дульной скорости в четвёртой степени на вероятность иметь снаряд в руках ограничено сверху константой изделия,
\[
v^{4}\,P_{\text{возв}}\le v_{\text{кр}}^{4}=\mathrm{const}.
\]
Всякое улучшение изделия ухудшает возможность его исследования.
\end{corollary}

Следствие~\ref{cor:uncert} представляется авторам наиболее фундаментальным результатом работы. Оно объясняет, в частности, почему испытания~3 и~5, давшие наибольшую информацию о характеристиках изделия, не дали ни одного вещественного доказательства этих характеристик, и почему группа перешла от поисков снаряда к изучению мишеней~--- методологический переход, который в рамках наивного подхода выглядел бы как капитуляция, а в рамках настоящей теории является единственно возможной стратегией.

\begin{figure}[h!]
\centering
\begin{tikzpicture}
\begin{axis}[
  width=11.2cm, height=6.0cm,
  xlabel={Дульная скорость $v$, м/с}, ylabel={$P_{\text{возв}}$},
  ymode=log, xmin=140, xmax=650, ymin=0.003, ymax=1.6,
  grid=both, grid style={gray!20},
  tick label style={font=\scriptsize}, label style={font=\small},
  legend style={font=\scriptsize, at={(0.98,0.97)}, anchor=north east, fill=\zhplbadgecolor{white}{zhplBg}}]
\addplot[thick, blue!70!black, domain=158.5:614.7, samples=300] {(158.5/x)^4};
\addlegendentry{$P_{\text{возв}}=(v_{\text{кр}}/v)^4$, $v_{\text{кр}}=158{,}5$ м/с}
\addplot[dashed, gray] coordinates {(614.8,0.003) (614.8,1.6)};
\addlegendentry{$v_\infty=614{,}8$ м/с}
\addplot[only marks, mark=*, mark size=2pt, red!70!black] coordinates {(158.5,1.0) (350.8,0.0417) (461.4,0.0139)};
\addlegendentry{испытания 1, 3, 5}
\node[font=\scriptsize, anchor=west] at (axis cs:170,1.20) {исп.~1};
\node[font=\scriptsize, anchor=west] at (axis cs:360,0.070) {исп.~3};
\node[font=\scriptsize, anchor=east] at (axis cs:452,0.0250) {исп.~5};
\end{axis}
\end{tikzpicture}
\caption{Вероятность возврата снаряда как функция дульной скорости в логарифмическом масштабе. Наклон прямой равен~$-4$ и определяется исключительно тем, что дальность растёт как $v^2$, а площадь~--- как квадрат дальности; этот показатель не зависит от того, изотермично или адиабатично расширение газа в стволе. Точки испытаний нанесены по расчётным скоростям из~\eqref{eq:main}: исп.~1~--- выбитый запирающий элемент, исп.~3~--- заклинивание и выстрел вверх, исп.~5~--- штатный режим. Пунктиром отмечен недостижимый предел $v_\infty$ из теоремы~\ref{th:limit}: правее него ось скорости физически пуста.}
\label{fig:lost}
\end{figure}

\section{Психофизика оператора}

\subsection{Иллюзия отскока}

По итогам второго испытания у наблюдателей сформировалось устойчивое ложное впечатление, что снаряд \emph{отскочил} от мишени, не поразив её. Четвёртое испытание это впечатление разрушило: снаряд прошёл мишень насквозь. Мы утверждаем, что оба наблюдения были корректны с точки зрения зрительной системы и что ошибка носила не перцептивный, а байесовский характер.

Пусть $d$~--- толщина мишени (\SEC), $v$~--- скорость снаряда. Характерное время взаимодействия
\[
\tau=\frac{d}{v}\sim\frac{10^{-2}~\text{м}}{4{,}61\cdot10^{2}~\text{м/с}}\approx2{,}2\cdot10^{-5}~\text{с},
\]
тогда как порог временного разрешения зрительной системы человека $\tau_{\text{зр}}\approx4\cdot10^{-2}$~с. Отношение $\tau_{\text{зр}}/\tau\approx1{,}8\cdot10^{3}$.

\begin{theorem}[О ложном отскоке]\label{th:rebound}
Если $\tau\ll\tau_{\text{зр}}$, то гипотезы «снаряд пробил мишень» и «снаряд отскочил от мишени» наблюдательно неразличимы: в обоих случаях наблюдатель регистрирует единственное событие~--- совпадение снаряда с мишенью в одном кадре восприятия с последующим отсутствием снаряда в поле зрения. Апостериорное предпочтение определяется исключительно априорным распределением наблюдателя, то есть его мнением об изделии до выстрела.
\end{theorem}

\begin{proof}
Правдоподобия обеих гипотез при данных, состоящих из одного кадра, равны между собой, поскольку кадр не содержит информации о событиях длительностью $\tau$. По формуле Байеса апостериорное отношение равно априорному.
\end{proof}

Теорема~\ref{th:rebound} имеет неожиданное методологическое следствие: до испытания~4 группа считала изделие слабым, поэтому \emph{видела} отскок; после испытания~4, в котором мишень была пробита навылет и продукт пробития оказался доступен непосредственному осмотру, апостериорное распределение сместилось, и все предыдущие наблюдения были ретроспективно переинтерпретированы. Это единственный зафиксированный в настоящей работе случай, когда новые данные изменили мнение исследователей, и авторы полагают его достойным отдельного упоминания.

\subsection{Лексическая термометрия}

Наиболее устойчиво регистрируемым параметром испытаний оказалась не скорость снаряда и не глубина пробития, а плотность ненормативной лексики в речи оператора. Введём

\begin{definition}[Лексическая плотность]
Лексической плотностью $\lambda$ называется число ненормативных лексем, произнесённых оператором в единицу времени в первые шестьдесят секунд после выстрела, измеряемое в мин$^{-1}$.
\end{definition}

\begin{theorem}[Закон лексической термометрии]\label{th:lex}
Для испытаний, завершившихся поражением мишени, лексическая плотность связана с удельной энергией $E$, переданной мишени, соотношением
\[
\lambda=\lambda_0+k\,\ln\!\Bigl(1+\frac{E}{E_0}\Bigr),
\]
где $\lambda_0$~--- фоновая лексическая плотность оператора, $k$~--- лексическая восприимчивость, $E_0$~--- порог удивления.
\end{theorem}

Метод наименьших квадратов по четырём результативным испытаниям при $E_0=0{,}5$ (усл.\,ед.) даёт
\[
\lambda_0=2{,}09~\text{мин}^{-1},\qquad k=8{,}88~\text{мин}^{-1},\qquad R^2=0{,}997 .
\]
Испытание~3 (заклинивание, выстрел в воздух) даёт $\lambda=15{,}8$~мин$^{-1}$ при $E\approx0$ и является выбросом, отклоняющимся от модели более чем на $13$~мин$^{-1}$. Мы объясняем это тем, что ненормативная лексика оператора имеет две компоненты различной природы~--- восторженную, описываемую теоремой~\ref{th:lex}, и фрустрационную, зависящую не от энергии на мишени, а от её отсутствия. Построение двухкомпонентной модели выходит за рамки настоящей работы и, по мнению рецензента~3, за рамки жанра научной публикации вообще.

\begin{figure}[h!]
\centering
\begin{tikzpicture}
\begin{axis}[
  width=11.2cm, height=6.0cm,
  xlabel={Удельная энергия на мишени $E$, усл.\,ед.},
  ylabel={Лексическая плотность $\lambda$, мин$^{-1}$},
  xmin=0, xmax=4.0, ymin=0, ymax=30,
  grid=both, grid style={gray!20},
  tick label style={font=\scriptsize}, label style={font=\small},
  legend style={font=\scriptsize, at={(0.98,0.06)}, anchor=south east, fill=\zhplbadgecolor{white}{zhplBg}, cells={anchor=west}}]
\addplot[thick, blue!70!black, domain=0:4, samples=300] {2.09+8.88*ln(1+x/0.5)};
\addlegendentry{$\lambda=2{,}09+8{,}88\ln(1+E/0{,}5)$, $R^2=0{,}997$}
\addplot[only marks, mark=*, mark size=2pt] coordinates {(0.35,6.9) (1.0,11.5) (1.9,16.4) (3.4,20.2)};
\addlegendentry{испытания 1, 2, 4, 5}
\addplot[only marks, mark=x, mark size=3.5pt, thick, red!70!black] coordinates {(0.02,15.8)};
\addlegendentry{испытание 3 (выброс, фрустрационная компонента)}
\end{axis}
\end{tikzpicture}
\caption{Зависимость лексической плотности оператора от удельной энергии, переданной мишени. Кривая~--- результат метода наименьших квадратов по четырём точкам, а не проведённая от руки линия тренда. Испытание~3 показано отдельным маркером как принадлежащее другой компоненте.}
\label{fig:lex}
\end{figure}

\section{Сводка испытаний и вербальные протоколы}

\subsection{Сводка}

\begin{table}[H]
\centering
\small
\caption{Сводка пяти натурных испытаний изделия PWA.}
\footnotesize
\begin{tabular}{@{}clllc@{}}
\toprule
№ & Режим & Запирающий элемент & Исход & Снаряд \\
\midrule
1 & стенд & исходный        & элемент выбит, выстрела нет  & не найден \\
2 & стенд & \SECs{} (в.\,2) & поражение мишени             & возвращён мишенью \\
3 & с рук & \SECs{} (в.\,2) & заклинивание, выстрел вверх  & не найден \\
4 & с рук & \SECs{} (в.\,2) & сквозное пробитие мишени     & не найден \\
5 & с рук & \SECs{} (в.\,2) & пробитие мишени из \SECs     & не найден \\
\bottomrule
\end{tabular}
\normalsize
\end{table}

\subsection{Вербальные протоколы}

Ниже приводятся выдержки из аудиодорожки испытания~5, отобранные по критерию информативности. Транскрипция сохранена с минимальной редакторской правкой; часть лексем приведена в усечённом виде в соответствии с требованиями редакции. Авторы отмечают, что в исходном материале плотность подобных лексем существенно выше приводимой, и полная транскрипция была бы неотличима от результата, а не от его описания.

\noindent\fcolorbox{black}{gray!6}{\begin{minipage}{0.94\linewidth}
\small
\textbf{П-5.1} (момент пробития): «ух еба».\\[2pt]
\textbf{П-5.2} (обращение к оператору стенда): «а ты говоришь слабая \SEC».\\[2pt]
\textbf{П-5.3} (демонстрация мишени): «бл*, мы тут думали, что х*йня, но ни*уя, она пробила это, пробив вот эту \SEC{}\ldots»\\[2pt]
\textbf{П-5.4} (заключительное): «вы понимаете, что мы сделали? Я всё, ну реально, я в а*уе. С этим реально можно идти и \SEC».
\end{minipage}}

\vspace{2mm}
Протокол П-5.1 представляет интерес как реакция с латентностью менее $\tau_{\text{зр}}$, то есть предшествующая осознанному распознаванию события; протокол П-5.3 содержит единственную в корпусе явную формулировку апостериорного пересмотра гипотезы (теорема~\ref{th:rebound}), а протокол П-5.4 фиксирует переход от исследовательской постановки задачи к прикладной, содержание которой остаётся \SEC{} и обсуждению в настоящей работе не подлежит.

\begin{table}[H]
\centering
\small
\caption{Соответствие вербальных реакций количественным характеристикам испытаний.}
\begin{tabular}{@{}clcc@{}}
\toprule
№ & Характер реакции & $E$, усл.\,ед. & $\lambda$, мин$^{-1}$ \\
\midrule
1 & недоумение, осмотр установки              & 0,35 & 6,9 \\
2 & сдержанное одобрение, сомнение в пробитии & 1,00 & 11,5 \\
3 & фрустрация, поиск снаряда в траве         & $\approx0$ & 15,8 \\
4 & подтверждение сквозного пробития          & 1,90 & 16,4 \\
5 & полный пересмотр модели изделия           & 3,40 & 20,2 \\
\bottomrule
\end{tabular}
\end{table}

\section{Обсуждение и ограничения}

Настоящая работа обязана честно перечислить причины, по которым её результаты не следует использовать.

\textbf{1. Все числа внутренне согласованы и внешне не проверены.} Единственной первичной величиной, измеренной в эксперименте, является зависимость $v(x)$; все остальные параметры ($v_\infty$, $x_0$, $\gamma$, $\pi^\ast$, $v_{\text{кр}}$, $\lambda_0$, $k$) получены подгонкой к ней и друг к другу. Независимое хронографирование не проводилось. Адиабатическая модель описывает данные с $R^2>0{,}998$, но при $n=5$ испытаниях столь высокое согласие является скорее свойством метода наименьших квадратов, чем свойством природы.

\textbf{2. Информативные испытания уничтожают собственные данные.} Согласно следствию~\ref{cor:uncert}, единственный снаряд, оказавшийся пригодным для изучения, был получен в испытании с наименьшей переданной энергией. Все выводы о бронепробитии сделаны по состоянию мишени, то есть по объекту, который в момент измерения перестал быть тем объектом, свойства которого измерялись. Это ограничение неустранимо в принципе и не может быть преодолено увеличением выборки: каждое новое испытание уменьшает вероятность возврата снаряда, а не увеличивает.

\textbf{3. Три из пяти испытаний проводились с рук.} Теорема~\ref{th:plug} утверждает существование оптимума $\pi^\ast$, но ничего не утверждает о том, находилось ли фактическое значение $\pi$ вблизи него в каждом конкретном пуске. Испытания~1 и~3 показывают, что не находилось. При том же распределении отказов и ином его исходе на левой ветви кривой рис.~\ref{fig:plug} цена эксперимента измерялась бы не в утраченных снарядах, а в конечностях группы, и текст статьи диктовался бы, а не набирался. Авторы считают необходимым зафиксировать: изделие PWA является несертифицированным сосудом, работающим под давлением, находящимся вне какого-либо контроля, и удержание такого сосуда руками в момент срабатывания не является методикой, а является статистикой. Ни один из приведённых в работе результатов не оправдывает воспроизведения описанной постановки эксперимента, а сама постановка приводится здесь исключительно потому, что она уже состоялась.

\textbf{4. Лексическая термометрия не является измерением.} Теорема~\ref{th:lex} устанавливает корреляцию между $\lambda$ и $E$, однако $E$ в свою очередь оценивалась по состоянию мишени, интерпретация которого производилась теми же операторами, чья лексика измерялась. Корреляция, таким образом, неустранима и не может быть очищена от наблюдателя ни при какой организации эксперимента.

Совокупность пунктов 1--4 позволяет заключить: единственным строго установленным результатом настоящей работы является утверждение о том, что чем лучше работает изделие PWA, тем меньше о нём можно узнать. Этот результат, в отличие от остальных, не зависит ни от одного подгоночного параметра и потому может считаться единственным подлинно научным содержанием статьи.

\section{Заключение}

Изделие PWA, разработанное по заказу \SEC{} и рассекреченное в текущем году, впервые получает связное теоретическое описание. Настоящий текст был написан непосредственно по итогам испытаний, то есть в год их проведения, однако до снятия грифа публикации не подлежал и пролежал в закрытом фонде целиком; к печати он допущен только сейчас и только после необходимой цензурной обработки, в ходе которой из него были изъяты все конструктивные параметры изделия, наименование заказчика, состав мишеней и содержание прикладной части программы. Изъятия отмечены в тексте как \SEC{} и восстановлению читателем не подлежат; авторы сохранили формулировки в том виде, в каком они были записаны по горячим следам, включая те, которые в спокойной обстановке они сформулировали бы иначе. Показано, что его внутренняя баллистика адиабатична и выходит на строгий предел скорости, что запирающий элемент имеет единственный оптимум, что снаряд принципиально невозвратен, а иллюзия отскока имеет байесовскую природу. Установлен закон лексической термометрии, позволяющий оценивать бронепробитие по аудиодорожке с точностью, не уступающей точности прямых измерений, поскольку прямых измерений не производилось.

Полученные результаты, вместе с зарегистрированной глубиной пробития мишени из \SEC, представляются авторам достаточным основанием для дальнейшего финансирования, а также для расширения программы испытаний на стендовый режим, к которому группа настоятельно рекомендует вернуться.

\section*{Вклад авторов}
\textbf{Р.\,Г.~Александров}~--- концептуализация, методология, полевые испытания, написание черновика, административное сопровождение процедуры рассекречивания.
\textbf{Ю.\,П.~Михаилов}~--- полевые испытания, запись видеоматериала, транскрипция вербальных протоколов, а также единственный участник проекта, находившийся на линии огня во всех трёх пусках с рук.
\textbf{К.\,С.~Антропикович}~--- программное обеспечение, визуализация, формальный анализ, вычисление касательных в замкнутом виде, отказ от вычисления величин, помеченных \SEC, техническая поддержка бесконечного терпения.

\section*{Финансирование}
Работа выполнена при поддержке гранта Института Прикладной Метафизики № 2026-$\Omega$-$\infty$ «Разработка методов исследования объектов, исчезающих в момент исследования», а также в рамках закрытой темы, шифр и объём финансирования которой \SEC.

\section*{Доступность данных}
Первичные данные доступны по обоснованному запросу~--- а также по необоснованному, с той же вероятностью успеха. Метаемые тела, использованные в испытаниях~1, 3, 4 и~5, доступны на площадке проведения испытаний в объёме, определяемом теоремой~\ref{th:lost}. Видеоматериалы испытаний прилагаются к настоящей публикации в качестве иллюстративного приложения и научной ценности, превышающей их развлекательную, не имеют.

\begin{thebibliography}{9}
\bibitem{adiabatov} \emph{Адиабатов И.\,И.} О расширении газов в каналах, о которых лучше не спрашивать. Труды Института Прикладной Метафизики, б.г.
\bibitem{zapirailo} \emph{Запирайло П.\,В., Отпиркин С.\,С.} Двойственность запирающего элемента: между потерей снаряда и потерей оператора. М.: Изд-во ИПМ, 2025.
\bibitem{nevozvratov} \emph{Невозвратов А.\,Ф.} Теория поиска предметов, которые уже никогда: монография. 2-е изд., дополненное сведениями о том, что искать бесполезно. 2026.
\bibitem{leksikonov} \emph{Лексиконов Б.\,Б.} Ненормативная лексика как первичный измерительный прибор. Вестник прикладной лженауки, 2026, № 1.
\end{thebibliography}

\noindent\textbf{Конфликт интересов.} Не заявлен, хотя, согласно следствию~\ref{cor:uncert}, строго доказать его отсутствие невозможно: всякая процедура проверки конфликта интересов уменьшает вероятность сохранения предмета проверки в состоянии, пригодном для проверки.

\bigskip
\noindent\hrule
\smallskip
\noindent\footnotesize\emph{Примечание редакции.} Рукопись поступила в редакцию в год проведения испытаний и всё это время хранилась в закрытом фонде без права публикации. Гриф снят постановлением №\,47-$\Omega$ от 06.09.2026; текст печатается с сокращениями, внесёнными в порядке обязательной цензурной обработки. Редакция обращает внимание читателей на то, что все конструктивные параметры изделия PWA сохраняют гриф \SEC{} и в настоящей публикации намеренно не восстановимы. Редакция не рассматривает статью как руководство, инструкцию или приглашение и категорически не рекомендует воспроизведение описанных испытаний, в особенности в режиме удержания изделия руками.

\end{document}
